\section{Classical global inputs with fixed doubling normalization}
\label{sec:sr-classical-doubling-inputs}

We match the two classical global inputs used by the theta
comparison to their precise published statements.
This establishes the anisotropic Siegel--Weil and Eisenstein
continuation inputs themselves.
It does not replace the separate local quotient comparison,
rank-recursion and covolume arguments.

Let \(E/F\) be a quadratic CM extension, with involution
\(c\), and fix a nontrivial character
\(\psi:\mathbb A_F/F\to\mathbb C^\times\).
Use the source's choice
\(\psi_\infty(t)=\exp(-2\pi i\operatorname{Tr}_{F_\infty/\mathbb R}t)\)
and the selfdual measures for it.
Let \(Y\) be a nondegenerate Hermitian or skew-Hermitian
space of dimension \(n\), and put
\[
             Y^\square=Y\oplus Y^{-},\qquad
             G^\square=U(Y^\square).
\]
The diagonal \(Y^\triangle=\{(y,y):y\in Y\}\) is a rational
maximal isotropic subspace.
Its stabilizer \(P\) is the Siegel parabolic, with Levi
\(\operatorname{Res}_{E/F}\operatorname{GL}_n\).
Thus the doubled group is quasi-split even when \(U(Y)\)
is an inner form of a different signature.

Write \(|\cdot|_E\) for the normalized adelic absolute value,
and let \(K\) be a standard maximal compact subgroup in
a rational hyperbolic frame carrying \(Y^\triangle\) to the
first isotropic coordinate space.
For the Iwasawa height
\(a(m(t)uk)=|\det t|_E\), use normalized induction
\begin{equation}\label{eq:sr-exact-induced-section}
 f_s(m(t)ug)
       =\chi(\det t)|\det t|_E^{\,s+n/2}f_s(g).
\end{equation}
The shift \(n/2\) is forced by the parabolic modulus.
Over a splitting field the unipotent radical is the space of
Hermitian \(n\)-by-\(n\) matrices.
For a diagonal \(t=\operatorname{diag}(t_i)\), its diagonal
coordinates contribute \(\prod_i|t_i|_E\), and each
off-diagonal \(E\)-coordinate contributes
\(|t_i|_E|t_j|_E\).
Their product is \(|\det t|_E^n\).
Conjugation invariance and triangular decomposition give the
same formula for every \(t\).
Hence \(\delta_P^{1/2}=|\det t|_E^{n/2}\).

\subsection{The definite theta integral and its constant}

Both ranks satisfy \(m,n\ge1\): these are explicit hypotheses
in the source's \srcid{SiegelWeilStd}, not restrictions added
for the application.  No zero-rank case is quantified there.
Let \(A\) have dimension \(m\), be of the opposite
Hermitian parity to \(Y\), and be definite at every infinite
place of \(F\).
Choose the standard unitary splitting character
\[
                    \chi|_{\mathbb A_F^\times}
                                      =\eta_{E/F}^{\,m}.
\]
Let \(\omega^\square\) be the corresponding Weil representation
for \(U(A)\times G^\square\) on its adelic Schwartz space.
For a Schwartz vector \(\Phi\) that is \(K\)-finite, the
Siegel--Weil section is
\[
 s_0=(m-n)/2,\qquad
 f_{\Phi,s}(g)
     =a(g)^{\,s-s_0}
                  (\omega^\square(g)\Phi)(0).
\]
The Weil action on the Siegel Levi has exponent \(m/2\).
Since \(s_0+n/2=m/2\), this is exactly
\eqref{eq:sr-exact-induced-section}, rather than an induction
shift chosen after the identity.

We use the following published theorem: the anisotropic
Siegel--Weil formula is holomorphic at \(s_0\) and equals
\(c_{m,n}\) times the probability-normalized theta integral,
where \(c_{m,n}=1\) for \(m>n\) and \(c_{m,n}=2\) for
\(m\le n\).  This is Theorem~7.1 of Gan--Qiu--Takeda,
\emph{The regularized Siegel--Weil formula (the second term
identity) and the Rallis inner product formula},
\url{https://arxiv.org/abs/1207.4709}.

Here is the full normalization check for our application.
A nonzero \(F\)-isotropic vector in \(A\) would remain
isotropic at every infinite place, contradicting definiteness.
Thus \(U(A)\) is anisotropic over \(F\) and its adelic
quotient is compact.
For any Haar measure \(\mu\) and any measurable fundamental
domain \(D_A\), put
\[
       {\rm Av}_{A,\mu}(H)
           =\mu(D_A)^{-1}\int_{D_A}H(h)\,d\mu(h).
\]
Its volume is finite and positive.
Replacing \(\mu\) by \(t\mu\), \(t>0\), multiplies numerator
and denominator by the same \(t\); changing fundamental
domain preserves the integral of an automorphic function.
The normalized integral therefore agrees with the
Tamagawa-normalized integral of the theorem.
For unitary groups the extra displayed factor
\(\tau(U(A))/[E:F]\) in that theorem is \(1\).
Consequently its conclusion is precisely
\begin{equation}\label{eq:sr-exact-anisotropic-sw}
 E(g,s_0,f_{\Phi})
       =\begin{cases}1,&n<m,\\2,&m\le n\end{cases}
          {\rm Av}_{A,\mu}
            \left(\sum_{x}
                  \omega^\square(g,h)\Phi(x)\right).
\end{equation}
All finite components of \(\Phi\), including ramified ones,
are arbitrary Schwartz vectors in this theorem.
No spherical-vector restriction or later adjustment of the
ramified Haar measures occurs in
\eqref{eq:sr-exact-anisotropic-sw}.
For the other Hermitian parity, choose
\(b\in E^\times\) with \(\overline b=-b\).
Replace the form on \(A\) by \(b\) times that form and the
form on \(Y\) by \(b^{-1}\) times that form.
Their Hermitian parities are reversed, their unitary groups
are unchanged, and their tensor symplectic pairing is
unchanged because the two scalars multiply to one.
This gives the skew-definite version with exactly the same
\(\psi\), measures and factor.

The source permits a rational Darboux realization \(d\), a
Siegel evaluation functional \(l\), and an action \(S\) of
\(U(A)\), rather than prescribing coordinate evaluation
in advance.
Its \srcid{IsDbl} conditions provide exactly their
normalization: rational vectors have rational Darboux
coordinates; \(l\) is invariant under the diagonal
modulations and under \(S\); and
\[
 \sum_{y\in A\otimes Y}
      l\bigl(\rho_d(y,-y)\Phi\bigr)
                               =\sum_x\Phi(x).
\]
There is a direct scalar check, requiring no unproved
uniqueness assertion for distributions.
Choose a rational polarization carrying the diagonal
\(A\otimes Y^\triangle\) to the modulation space.
Its opposite Lagrangian is \(A\otimes\{(y,-y)\}\).
The source's rationality condition makes its rational
translations the full \(F\)-lattice of that opposite
Lagrangian.  Rational changes of Schr\"odinger polarization
preserve the theta functional: linear changes reindex the
rational sum and Fourier changes use Poisson summation
with selfdual measures and the product formula.
Thus in these coordinates, writing
\[
 T_r\phi(t)=\phi(t+r),\qquad
 M_b\phi(t)=\psi(b\cdot t)\phi(t),
\]
the two conditions on \(l\) become the following statements
for \emph{all adelic} \(b\) and \emph{every} Schwartz \(\phi\):
\[
 l(M_b\phi)=l(\phi),\qquad
       \sum_{r\in F^{2mn}}l(T_r\phi)
                         =\sum_{r\in F^{2mn}}\phi(r).
\]
The rational Darboux change identifies the opposite
translation lattice with \(F^{2mn}\).
Apply the second equality also to \(M_b\phi\).
Since \(T_rM_b=\psi(b\cdot r)M_bT_r\), the first equality
gives
\[
 \sum_{r\in F^{2mn}}\psi(b\cdot r)
                 \bigl(l(T_r\phi)-\phi(r)\bigr)=0
                                     \quad(b\in\mathbb A_F^{2mn}).
\]
The series is absolutely summable: the source states an
unconditional complex \srcid{HasSum} for the left side at
\(b=0\), and the Schwartz theta sum on the right is
absolutely convergent.
Multiplication by the unitary characters gives one uniform
absolute majorant for every \(b\).
Integrating over the compact group
\((\mathbb A_F/F)^{2mn}\), of mass one, and using character
orthogonality selects \(r=0\).
It gives \(l(\phi)=\phi(0)\) for every Schwartz \(\phi\).
Thus both the evaluation functional and its scalar are
fixed by the exact \srcid{IsDbl} conditions.
In particular no freedom to multiply \(l\) by a nontrivial
constant remains.

The unspecified \(U(A)\)-action is fixed as well.
In the diagonal polarization \(U(A)\) preserves both
Lagrangians and acts on the translation coordinate space
\(X=(A\otimes_EY)(\mathbb A_F)\) by a linear map \(a_h\).
Its \(F\)-determinant has absolute value one:
the unitary determinant identity gives
\(N_{E/F}(\det_E h)=1\), and the determinant on \(X\)
is its \(n\)-th power.
Consequently the ordinary Weil action on this factor is
\[
                    (R_h\phi)(x)=\phi(a_h^{-1}x);
\]
there is no Jacobian factor.  It preserves Schwartz
functions, fixes evaluation at zero and satisfies the
same Heisenberg intertwining relation as \(S(h)\).
We now use those facts, rather than presume scalar
uniqueness of a representation.
For every Schwartz \(\phi\) and every Heisenberg
parameter \(w\), the intertwining relation in
\srcid{IsDbl.lin} and the equality \srcid{IsDbl.fix}
give
\[
 \begin{split}
 l\!\left(\rho(w)S(h)\phi\right)
   &=l\!\left(S(h)\rho(h^{-1}w)\phi\right)\\
   &=l\!\left(\rho(h^{-1}w)\phi\right)
    =l\!\left(\rho(w)R_h\phi\right).
 \end{split}
\]
The same relation holds in the transported rational
Darboux realization.
Since \(l=\operatorname{ev}_0\), taking \(w=(x,0)\)
among the pure translations says
\((S(h)\phi)(x)=(R_h\phi)(x)\) for every \(x\).
Thus \(S(h)=R_h\) on the entire Schwartz space.
This proves the compatibility with the ordinary
\(U(A)\)-Weil action, including its scalar, using precisely
the preservation, Egorov and fixed-functional conditions
present in \srcid{IsDbl}.  Behavior on functions outside
the Schwartz space is irrelevant to either side of the
identity.

We also check that the unspecified doubled action
\(\mathrm{om}\) has no remaining scalar character.
The exact definition of \srcid{IsUWeilSplit} is
\[
  \srcid{IsLagr}\ \wedge\
  \srcid{IsUSplitting}\ \wedge\ \srcid{SiegelRule}.
\]
Here \srcid{IsUSplitting} includes preservation of the
Schwartz space, Egorov intertwining, preservation of
the \(L^2\)-norm for Haar measure, and continuity after
every admissible linear functional.
Norm preservation and the inverse group element extend
each operator to a unitary operator on \(L^2(X)\).
The stated continuity gives continuous matrix
coefficients against Schwartz vectors; these are
admissible functionals, since pairing with a fixed
Schwartz vector is continuous on each fixed finite
Schwartz factor.
Density and unitarity give a continuous unitary action.

Compare this action with the ordinary Weil action for
the specified \(\chi\).
For each \(g\), their ratio commutes with all adelic
translations and modulations, by Egorov.
Such a bounded operator is scalar: commuting with
modulations makes it multiplication by a bounded
measurable function, and commuting also with
translations makes that function constant almost
everywhere.
The first assertion follows by Fourier inversion for
compactly supported test functions and then the
bounded multiplication spectral algebra.
Thus the ratio is \(c(g)\in U(1)\), and the group law
and continuity make \(c\) a continuous character.
Apply \srcid{SiegelRule} to
\(l=\operatorname{ev}_0\), which is admissible and
invariant under the diagonal modulation space.
Since the ordinary action obeys the same rule, a
Schwartz vector with \(\phi(0)\ne0\) gives \(c(p)=1\)
for every element of the adelic Siegel parabolic.

At each local place, a character trivial on this
parabolic is trivial on the whole doubled group.
Indeed the parabolic contains a Borel torus and its
positive unipotent root groups; characters are
conjugation invariant, so they are trivial on the
opposite root groups as well.
The root-group expressions for the Weyl
representatives and Bruhat decomposition then
generate the local group.
This applies also at split places, where the group
is general linear.
Finite products of local factors are dense in the
adelic restricted product, so continuity gives
\(c(g)=1\) everywhere.
The two actions agree on \(L^2\), and on Schwartz
functions agree pointwise by continuity and full
support of Haar measure.

The splitting character and doubled action in
\srcid{IsDbl.weil} are therefore exactly those in the
classical formula, not a representation identified
only up to an unspecified scalar.
Finally \srcid{IsHgt.val} gives the height just used.
The finite-dimensional \(K\)-orbit condition is precisely
\(K\)-finiteness.
Thus \eqref{eq:sr-exact-anisotropic-sw}, including its
holomorphy assertion, establishes the classical input
\srcid{SiegelWeilStd} with the source's quantifiers.
The later change from a standard hyperbolic frame to a
lattice-adapted frame transports this identity; it does not
change the factor \(1\) or \(2\).

\subsection{Meromorphic continuation in the required growth space}

For a continuous \(K\)-finite section \(f=f_{s_1}\) obeying
\eqref{eq:sr-exact-induced-section}, set
\(f_s=a^{s-s_1}f\) and initially
\[
                  E(g,s,f)
                     =\sum_{\gamma\in P(F)\backslash
                                  G^\square(F)}f_s(\gamma g).
\]
This quotient agrees with the rational parabolic quotient
in the source, since every rational maximal isotropic
frame is obtained from the fixed diagonal one.

We use Bernstein--Lapid, \emph{On the meromorphic
continuation of Eisenstein series}, JAMS~37 (2024), 187--234,
Theorem~2.3 and the corollary at the end of Section~8,
\url{https://doi.org/10.1090/jams/1020}.
They give convergence in a right half-plane and
meromorphic continuation; locally, a scalar polynomial
clears the poles into one Fr\'echet space of functions of
uniform moderate growth.

We check the section and topology hypotheses.
Restriction of \(f\) to \(K\) belongs to a finite-dimensional
continuous representation.
At the infinite places its matrix coefficients are smooth.
At finite places its compact image in a complex general
linear group is a totally disconnected compact Lie group,
hence finite; the section therefore has an open compact
stabilizer.
The parabolic transformation law extends this smoothness
from \(K\) to \(G^\square\).
The inducing representation of the Levi is the automorphic
character \(\chi\circ\det\), with the displayed spectral
twist.
It is an automorphic form, so the theorem applies to this
inducing datum without a cuspidality assumption.

At a fixed \(s_*\), restrict the theorem's several-variable
parameter to the one Siegel line \(s\).
Shrink a disk around \(s_*\) so that no other singular
hyperplane meets it.
A product of linear factors through \(s_*\) restricts to a
nonzero constant times \((s-s_*)^k\).
The right half-plane of convergence ensures that the
Siegel line is not contained in a singular hyperplane.
Thus, for some \(k\ge0\),
\[
                 H(s)=(s-s_*)^kE(\,\cdot\,,s,f)
\]
extends holomorphically across \(s_*\) in one fixed
uniform-moderate-growth Fr\'echet space.
Its zeroth seminorm supplies a fixed integer \(N\) and,
on a smaller closed disk, a fixed \(C\) with
\[
                     |H(s)(g)|\le C\,{\rm ht}(g)^N .
\]
Evaluation at \(g\), and all continuous linear growth
functionals, are holomorphic on this disk.

For clarity the latter statement gives exactly the weak
functional condition used by the source.
If a linear functional \(\lambda\) obeys
\[
 |\lambda(F)|\le C_N
       \quad\hbox{whenever}\quad |F(g)|\le{\rm ht}(g)^N ,
\]
homogeneity gives
\(|\lambda(F)|\le C_N\sup_g|F(g)|{\rm ht}(g)^{-N}\).
It is therefore continuous on the fixed growth space in
which \(H\) is holomorphic.
The composition \(s\mapsto\lambda(H(s))\) is holomorphic.
The height in the source is
\({\rm ht}_2(g)={\rm heightGL}_F(g)\) for the displayed
faithful matrix realization.
It is polynomially comparable to the algebraic matrix
height in the continuation theorem.
Indeed coordinates of either fixed faithful realization
are regular functions in the coordinates and inverse
determinant of the other; bounding the finitely many
coefficients and applying the local maximum norm bounds
these functions by a constant times a fixed power.
At all but finitely many finite places the coefficients
are integral, so the product of these constants is finite.
Applying this argument in both directions gives the
comparison.  The field, rank and matrix realizations are
fixed in this application; its constants may depend on
them, as the quantifiers of \srcid{EisBL} permit.
Multiplying the height by a constant or taking an
equivalent matrix height only alters \(C,N\).
This proves the \srcid{EisBL} input, including its
\srcid{IsGrowForm} quantifier, rather than only pointwise
meromorphic continuation.

These two propositions close the named classical inputs
to \srcid{theoremK\_head2}.
They do not by themselves prove its complete output:
\srcid{ThetaIntegralsAgree} additionally contains the
normalized real and finite quotient maps, all matching
identities, nonzero ordinary integrals and both
codimension-two recursions with character-independent
constants.
Those consequences require their displayed local and
global comparison arguments.
Nor do the two classical inputs imply
\srcid{CovolumeBoundOfX}, which belongs to the separate
integral lattice and covering-tower branch.
